\documentstyle[%


righttag%


,11pt%


,amssymb%


,amscd%


,russian%               remove in Eng.


,izvru%                 Set izven% in Eng.


% ,izvru-ks             for brief communications


]{amsart}





%


\newcommand{\Alb}{\operatorname{Alb}}


\newcommand{\tts}[1]{}





%\hoffset=-15mm%                  For Eng.


\setcounter{page}{80}


\god{2001}


\nomer{7~(470)} %                        Remove in Eng.


%\nomer{Vol.\,45, No.\,7}%               Set in Eng.


%\str{\pageref{begin}--\pageref{end}}%  Set in Eng.


\udk{512.6}





\author{\it М.Г.\,Танкеева}


% NB:   ^^ remove in Eng.


\title{Об аналитических циклах на комплексных торах}





\date{}


%\pagestyle{myheadings}%         Set in Eng.


\pagestyle{plain} %              Remove in Eng.


\begin{document}


%\thispagestyle{plain}%          Set in Eng.


\maketitle


\label{begin}





Степень трансцендентности поля мероморфных функций компактного комплексного


многообразия является важнейшим инвариантом, связанным с геометрией


многообразия. Например, из хорошо известной классификации


компактных комплексных поверхностей следует, что


на поверхности с одной алгебраически независимой


мероморфной функцией существует пучок эллиптических кривых.


Еще более экзотические


многообразия без непостоянных мероморфных функций почти не изучены и


традиционно


рассматриваются в литературе при построении примеров, опровергающих


обычную геометрическую интуицию [1].





Целью этой работы является изучение геометрии простых комплексных торов


без непостоянных мероморфных функций, что позволяет получить некоторую


информацию о геометрии многообразий, близких к таким торам в смысле


теории деформаций комплексных структур.





\begin{thm}


Пусть $V$ --- связное гладкое замкнутое


комплексно-аналитическое подмногообразие комплексного тора $T$.


Если $0<\dim (V)<\dim (T)$, то $V$ нестабильно


относительно малых деформаций тора $T$. Более того, если $N$ --- нормальное


расслоение к $V$ в $T$, то $H^1(V,N)\neq(0)$.


\end{thm}





Известно ([2], теорема 1), что из равенства $H^1(V,N)=(0)$


следует стабильность.


Пусть $X\to S$ --- версальное семейство Кураниши деформаций тора


$T=X_0$, $\Delta\subset S$ --- множество таких точек $s\in S$, что тор


$X_s$ прост и на $X_s$ нет непостоянных мероморфных функций. Хорошо


известно, что $\Delta$ плотно в $S$.





\begin{thm}


На простом комплексном торе $T$ без непостоянных


мероморфных функций любой аналитический цикл коразмерности не меньше $1$


является $0$-мерным.


\end{thm}





\begin{pf}


Известно [1], что на $T$ нет ненулевых дивизоров. Докажем, что на $T$ нет


кривых. Пусть $C\subset T$ --- неприводимая кривая. Можно считать, что


$0\in C$. Положим $C_n=C+\dotsb+C$ ($n$ слагаемых). Если $C_{n-1}\neq T$,


то $\dim (C_n)=n$. Действительно, $C_{n-1}\subset C_n$ и


$C_{n-1}\neq C_n$ (иначе $C_{n-1}=C_{n-1}+C_{n-1}$ и, следовательно,


$C_{n-1}$ --- нетривиальная аналитическая подгруппа в $T$, что противоречит


простоте тора $T$), поэтому $\dim (C_n)=1+\dim (C_{n-1})$,


так что $\dim (C_n)=n$. Если $g=\dim (T)$, то $C_{g-1}$ ---


дивизор на $T$, что противоречит теореме В.А.\,Краснова [1].





Пусть $S\subset T$ --- неприводимое приведенное замкнутое комплексное


подпространство минимальной строго положительной размерности на $T$.


Ясно, что $\dim (S)>1$, $\text{Sing}\,(S)$ конечно,


на $S$ нет нетривиальных дивизоров и поэтому нет непостоянных мероморфных


функций. Разрешая (изолированные) особенности $S$ с помощью


последовательных моноидальных преобразований $T$ с неособыми центрами ([3],


глава 0, \S\,7, теорема $I'$), получим гладкую модель $S'$ пространства $S$,


причем $S'\subset T'$, где $T'$ получается из кэлерова многообразия


$T$ с помощью последовательных раздутий неособых центров, а поэтому


$T'$ --- кэлерово многообразие. Следовательно, $S'$ --- кэлерово


многообразие без непостоянных мероморфных функций. Для таких многообразий


отображение Альбанезе $\alpha:S'\to\Alb (S')$ сюръективно [1].


В силу универсального свойства отображения Альбанезе существует


коммутативная диаграмма


$$


\begin{CD}


S' @> \alpha >> \Alb (S')\\


@VV \tau V @VV h V\\


S @> i >> T,


\end{CD}


$$


где $\tau:S'\to S$ --- бимероморфный морфизм, $i$ --- каноническое вложение,


$h$ --- композиция гомоморфизма торов и сдвига. Ясно, что


$h(\Alb (S'))$ --- подтор в $T$, содержащий $i(S)$, поэтому из


простоты $T$ следует сюръективность $h$. Из сюръективности $\alpha$ и $h$


следует неравенство $\dim (S')\geq\dim (T)$. Поэтому


$S=T$.


\end{pf}





\begin{thebibliography}{9}





\bibitem{1}


Краснов В.А. {\em Компактные комплексные многообразия без мероморфных функций}


// Матем. заметки. -- 1975. -- Т.\,17. -- \No\,1. -- C.\,119--122.


\bibitem{2}


Kodaira K. {\em On stability of compact submanifolds of complex manifolds} //


Amer. J. Math. -- 1963. -- V.\,85. -- P.\,79--94.


\bibitem{3}


Хиронака Х. {\em Разрешение особенностей алгебраических многообразий


над полями характеристики нуль} // Математика (сб. перев.). -- 1965.


-- Т.\,9. -- \No\,6. -- C.\,2--70.





\end{thebibliography}





\vskip 2cm





\post{\it Владимирский государственный университет}{\it Поступила}


\post{}{05.06.2000}





\label{end}


\end{document}





